\documentclass[12pt,a4paper]{article}

\usepackage[top=2.5cm, bottom=2.5cm, left=2.5cm, right=2.5cm, headheight=15pt]{geometry}
\usepackage{amsmath, amssymb}
\usepackage[T1]{fontenc}
\usepackage[utf8]{inputenc}
\usepackage{booktabs}
\usepackage{array}
\usepackage{tabularx}
\usepackage{longtable}
\usepackage{enumitem}
\usepackage{xcolor}
\usepackage{mdframed}
\usepackage{titlesec}
\usepackage{listings}
\usepackage{tcolorbox}
\usepackage{graphicx}
\usepackage{multirow}
\usepackage{hyperref}
\usepackage{parskip}
\usepackage{setspace}
\usepackage{float}
\usepackage{placeins}

\renewcommand{\familydefault}{\sfdefault}
\usepackage[scaled=1]{helvet}
\usepackage[helvet]{sfmath}

\tcbuselibrary{skins, breakable, listings}

\definecolor{headerblue}{RGB}{0, 56, 101}
\definecolor{sectionblue}{RGB}{30, 90, 160}
\definecolor{lightblue}{RGB}{220, 235, 252}
\definecolor{codebg}{RGB}{245, 247, 250}
\definecolor{codeframe}{RGB}{180, 200, 230}
\definecolor{answerbg}{RGB}{255, 253, 230}
\definecolor{answerframe}{RGB}{200, 180, 80}
\definecolor{notebg}{RGB}{240, 248, 240}
\definecolor{noteframe}{RGB}{100, 170, 100}
\definecolor{warningbg}{RGB}{255, 243, 220}
\definecolor{warningframe}{RGB}{220, 150, 50}
\definecolor{instrbg}{RGB}{245, 247, 250}
\definecolor{instrframe}{RGB}{140, 160, 190}
\definecolor{darkgray}{RGB}{60, 60, 60}

\definecolor{myColorA}{RGB}{3, 28, 46}%rightfootline
\definecolor{myColorB}{RGB}{13, 74, 119}%{5,0,70}%centerfootline
\definecolor{myColorC}{RGB}{255, 151, 24}%{0,0,100}%leftfootline

\titleformat{\section}
  {\large\bfseries\color{headerblue}}
  {\thesection.}{0.5em}{}
  [\color{sectionblue}\hrule height 0.8pt]

\titleformat{\subsection}
  {\normalsize\bfseries\color{sectionblue}}
  {\thesubsection}{0.5em}{}

\setlength{\parskip}{6pt}

\lstdefinestyle{pythonstyle}{
  language=Python,
  backgroundcolor=\color{codebg},
  basicstyle=\ttfamily\small,
  keywordstyle=\bfseries\color{sectionblue},
  commentstyle=\itshape\color{gray},
  stringstyle=\color{orange!70!black},
  showstringspaces=false,
  breaklines=true,
  frame=single,
  rulecolor=\color{codeframe},
  framesep=5pt,
  aboveskip=8pt,
  belowskip=4pt,
  xleftmargin=8pt,
  xrightmargin=4pt,
  numbers=left,
  numberstyle=\tiny\color{gray},
  numbersep=8pt,
}

\newtcolorbox{answerbox}[1][]{
  colback=answerbg, colframe=myColorC,
  title={\textbf{Answer space}},
  fonttitle=\small\bfseries,
  coltitle=black,
  breakable, #1
}

\newtcolorbox{notebox}[1][]{
  colback=notebg, colframe=noteframe,
  left=6pt, right=6pt, top=4pt, bottom=4pt,
  breakable, #1
}

\newtcolorbox{warningbox}[1][]{
  colback=warningbg, colframe=warningframe,
  left=6pt, right=6pt, top=4pt, bottom=4pt,
  breakable, #1
}

\newtcolorbox{instructionbox}{
  colback=lightblue, colframe=myColorB,
  left=8pt, right=8pt, top=6pt, bottom=6pt,
  breakable
}

\newtcolorbox{stepbox}[1][]{
  colback=instrbg, colframe=myColorA,
  title={\textbf{Instructions}},
  fonttitle=\small\bfseries,
  coltitle=white,
  breakable, #1
}

\newcommand{\answerline}{\vspace{0.9cm}\noindent\hrulefill\par}
\newcommand{\shortanswerlines}[1]{%
  \foreach \i in {1,...,#1}{\answerline}%
}

\graphicspath{{figures/}}
\newcommand{\figwide}[3]{\begin{figure}[htbp]\centering
\includegraphics[width=\textwidth]{#1}\caption{#2}\label{#3}\end{figure}}
\newcommand{\figpair}[4]{\begin{figure}[htbp]\centering
\includegraphics[width=0.62\textwidth]{#1}\hfill\includegraphics[width=0.33\textwidth]{#2}
\caption{#3}\label{#4}\end{figure}}

\begin{document}

\begin{center}
  {\LARGE\bfseries\color{headerblue}
    Coding Assignment 4\\[4pt]
    Convolutional Neural Networks}\\[10pt]
\end{center}

\section*{Overview}

This activity walks you through building, training, and evaluating convolutional neural network (CNN) architectures for multi-class medical image
classification using the \textbf{Brain Tumor MRI Dataset}. Beyond architecture
comparison, you will run six controlled experiments isolating the effects of
\textbf{loss function}, \textbf{optimizer}, \textbf{kernel size}, \textbf{stride},
\textbf{padding}, and \textbf{pooling strategy} on feature extraction and
classification performance. You will also produce, for \emph{every diagnostic class} in the
dataset, a worked demonstration of (a) the data augmentations applied to a
representative image and (b) a hand-verified convolution calculation showing
the input patch/image, kernel, receptive field, and resulting feature map.

\begin{instructionbox}
\textbf{Dataset:} Brain Tumor Magnetic Resonance Imaging\\
\textbf{Source:} \href{https://huggingface.co/datasets/alfa-bravo/Brain-Tumor-Mri-Dataset/blob/main/Brain\%20Tumor.rar}{Brain Tumor MRI Dataset (Hugging Face)}\\
\textbf{Classes:} 4
\textbf{Task:} Multi-class image classification, not binary (tumor/no-tumor).
\end{instructionbox}

\begin{warningbox}
\textbf{\textcolor{red}{Warning:}} Complete all coding tasks in your own Jupyter notebook,
following the instructions given in each part of this sheet. Answer the written
questions \emph{in this sheet} based on your own observations and course concepts.
\textbf{Do not copy code outputs verbatim as answers}, and \textbf{do not fabricate or
estimate experimental results} --- every number in your tables and every figure you
reference must come from a run you actually performed. Explain using your own words
and equations.
\end{warningbox}

\textbf{General instructions:}
\begin{itemize}[leftmargin=*]
  \item Implement all experiments in \textbf{Python} using \textbf{PyTorch}.
  \item Use \textbf{Google Colab (GPU runtime)} or a local machine with a CUDA-capable
        GPU where available. If you have no GPU access, reduce image resolution, batch
        size, and epoch count, and state the adjustment explicitly in your report.
  \item Organize your source code into reusable modules (Section~15). Do not write
        one monolithic script.
  \item Keep a fixed random seed (\texttt{SEED = 42}) for every experiment unless a
        question specifically asks you to vary it.
\end{itemize}

\section*{Learning Objectives}

By the end of this activity, you should be able to:
\begin{enumerate}[leftmargin=*, label=\arabic*.]
  \item Implement and compare CNN architectures for a multiclass medical image classification task.
  \item Design clinically-justified data augmentation strategies for MRI images and
        visually verify them on a representative sample from \emph{every} diagnostic class.
  \item Perform manual convolution calculations for a representative image of \emph{every}
        class and verify the computation against PyTorch.
  \item Compare multi-class loss functions appropriate for this task and
        relate their gradients to observed training behavior.
  \item Compare gradient-based optimizers appropriate for image classification and relate their
        update rules to observed convergence speed and generalization.
  \item Quantify how convolutional kernel size, stride, and padding affect output
        dimensions, receptive field, parameter count, and classification accuracy.
  \item Distinguish max pooling, average pooling, and global average pooling in terms
        of information retained, gradient flow, and generalization.
  \item Connect every experimental observation to the architectural or mathematical
        principle that explains it, supported only by results you actually obtained.
  \item Produce a reproducible, modular PyTorch codebase and a technical report
        documenting methodology, results, and interpretation.
\end{enumerate}

\section{Environment Setup and Data Exploration}

\subsection{Package Imports}

Copy and run the following cell at the top of your notebook. If any import fails,
install the missing package with \texttt{pip install <package>} before proceeding.

\begin{lstlisting}[style=pythonstyle]
import os, random, math, time
import numpy as np
import pandas as pd
import matplotlib.pyplot as plt
import seaborn as sns
from PIL import Image

import torch
import torch.nn as nn
import torch.nn.functional as F
from torch.utils.data import Dataset, DataLoader, random_split
import torchvision
from torchvision import transforms
from torchvision.datasets import ImageFolder
from sklearn.metrics import confusion_matrix, classification_report

SEED = 42
random.seed(SEED); np.random.seed(SEED); torch.manual_seed(SEED)

DEVICE = torch.device('cuda' if torch.cuda.is_available() else 'cpu')
print('Torch version :', torch.__version__)
print('Device        :', DEVICE)
\end{lstlisting}

\subsection{Download and Extract the Dataset}

The dataset is distributed as a single \texttt{.rar} archive on Hugging Face.

\begin{stepbox}
\begin{enumerate}[leftmargin=*, label=\arabic*.]
  \item Download \texttt{Brain Tumor.rar} from the dataset page linked in the
        Overview box above.
  \item Install an unrar backend (e.g.\ \texttt{pip install patool} together with
        the system \texttt{unrar}/\texttt{unar} utility, or \texttt{pip install rarfile}).
  \item Extract the archive into a local folder. Confirm that the extracted directory
        contains \textbf{four} subfolders, one per class: \texttt{glioma},
        \texttt{meningioma}, \texttt{pituitary}, \texttt{notumor} (folder names may
        vary slightly --- print them to confirm).
  \item Build a file list by walking the extracted directory, recording each image's
        file path and its class label taken from its parent folder name.
  \item Load the dataset with \texttt{torchvision.datasets.ImageFolder} (or a custom
        \texttt{Dataset} subclass) pointed at the extracted directory, and print the
        resulting \texttt{class\_to\_idx} mapping.
\end{enumerate}
\end{stepbox}

\begin{lstlisting}[style=pythonstyle]
import patoolib
patoolib.extract_archive('Brain Tumor.rar', outdir='brain_tumor_mri')

DATA_ROOT = 'brain_tumor_mri'   # update to the folder that contains the 4 class subfolders
full_dataset = ImageFolder(DATA_ROOT)
print('Classes found:', full_dataset.classes)
print('Class -> index map:', full_dataset.class_to_idx)
\end{lstlisting}

\subsection{Dataset Exploration}

\begin{stepbox}
\begin{enumerate}[leftmargin=*, label=\arabic*.]
  \item Count the number of images in each of the 4 class folders and display the
        counts as a table and as a bar chart.
  \item Display one representative image from each class in a single $1 \times 4$
        figure, with the class name as the subplot title. \textbf{Keep the index of
        each chosen representative image} --- you will reuse the \emph{same}
        representative image in Sections~3 and 4.
  \item Print the pixel dimensions and color mode (grayscale vs.\ RGB) of a handful
        of raw images from different classes.
\end{enumerate}
\end{stepbox}

\begin{tcolorbox}[colback=lightblue!40, colframe=myColorB, title={\textbf{Reflection Question 1}}, fonttitle=\small\bfseries, breakable]
From your class-count bar chart, is the dataset balanced across the four classes?
Report the exact counts you obtained. If the dataset is imbalanced, name
\textbf{two} concrete techniques (applicable later in training) that could
compensate for it.
\end{tcolorbox}

\begin{answerbox}
The dataset is not perfectly balanced: Glioma has 4,167 images (31.58\%),
Pituitary Tumor 3,232 (24.49\%), No Tumor 2,850 (21.60\%), and Meningioma
2,947 (22.33\%). Two ways to reduce majority-class bias are inverse-frequency
class weights in the loss and a class-balanced sampler or oversampling of the
less frequent classes. This notebook tested weighted cross-entropy and focal
loss with inverse-frequency weights.
\end{answerbox}

\begin{tcolorbox}[colback=lightblue!40, colframe=myColorB, title={\textbf{Reflection Question 2}}, fonttitle=\small\bfseries, breakable]
MRI images in this dataset may come from different scanners and therefore differ in
resolution and intensity range. Identify \textbf{two} preprocessing steps (besides
resizing) that address this, and explain, for each, what would go wrong if it were
skipped.
\end{tcolorbox}

\begin{answerbox}
First, convert every image to one grayscale channel. The raw sample included
grayscale PNGs and RGB JPEGs, so skipping this step would produce inconsistent
input channels and color representations. Second, scale pixel values to
$[0,1]$ and normalize with the training-split mean and standard deviation
($0.1726$ and $0.1837$). Without normalization, differences in scanner intensity
and image brightness can dominate the learned features; computing the statistics
from all splits would also leak information from validation and test data.
\end{answerbox}

\figwide{fig_class_counts.png}{Number of images per class (Reflection Question 1).}{fig:counts}
\figwide{fig_original_images.png}{One representative image per class, reused in Sections 3 and 4.}{fig:orig}
\FloatBarrier

\section{Data Preparation}

\begin{stepbox}
\begin{enumerate}[leftmargin=*, label=\arabic*.]
  \item Resize every image to a fixed size (e.g.\ $128 \times 128$ or $224 \times 224$,
        depending on your compute budget) and convert to a single channel (grayscale)
        or replicate to 3 channels if a pretrained architecture requires it.
  \item Split the dataset into \textbf{train / validation / test} using a
        \textbf{70\% / 15\% / 15\%} \emph{stratified} split (stratify by class label so
        the class proportions from Reflection Question~1 are preserved in every split).
        Fix \texttt{SEED = 42}.
  \item Compute the per-channel mean and standard deviation \textbf{from the training
        split only}, and use them to normalize all three splits.
  \item Wrap each split in a \texttt{DataLoader} with a consistent batch size
        (state your chosen value and justify it given your compute budget).
  \item Print the number of images and the class distribution in each of the three
        splits to confirm the stratified split worked as intended.
\end{enumerate}
\end{stepbox}

\begin{tcolorbox}[colback=lightblue!40, colframe=myColorB, title={\textbf{2.1 --- Questions}}, fonttitle=\small\bfseries, breakable]
\textbf{(a)} Why must normalization statistics be computed from the training split
only, and not the full dataset? \\[6pt]
\textbf{(b)} Why is a \emph{stratified} split preferable to a plain random split for
this dataset, given your answer to Reflection Question~1? \\[6pt]
\textbf{(c)} Holding the same splits, preprocessing, and random seed across every
architecture and every experiment in this assignment is essential. Explain the
controlled-experiment principle behind this requirement.
\end{tcolorbox}

\begin{answerbox}
\textbf{(a)} Using validation or test images to estimate normalization statistics
would leak information from those held-out sets into training and make evaluation
less independent. Training-only statistics let the validation and test sets act
as unseen data.\par
\textbf{(b)} Stratification keeps the class proportions similar across the
train, validation, and test splits. This matters because the class counts differ;
a plain random split could leave a smaller class under-represented in one split.\par
\textbf{(c)} Keeping the split, preprocessing, seed, and other settings fixed
controls confounding variables. The measured difference can then be attributed
more plausibly to the one factor being changed in each experiment.
\par
The run used grayscale $128\times128$ images, batch size 32, a 30-epoch maximum,
and patience 6. The split totals were 9,237 training, 1,979 validation, and 1,980
test images. Per-class train/validation/test counts were Glioma 2,917/625/625,
Meningioma 2,063/442/442, Pituitary Tumor 2,262/485/485, and No Tumor
1,995/427/428.
\end{answerbox}

\section{Data Augmentation --- Required for Every Class}

\begin{stepbox}
\begin{enumerate}[leftmargin=*, label=\arabic*.]
  \item Define a \texttt{torchvision.transforms} (or \texttt{albumentations}) pipeline
        containing \emph{clinically plausible} augmentations only:
        horizontal flip, small-angle rotation ($\pm 10$--$15^{\circ}$), brightness/contrast
        jitter, slight zoom or random-resized crop, and mild Gaussian noise.
  \item \textbf{Do not} use vertical flips or large elastic/shear distortions: these
        can produce anatomically implausible brain shapes that no longer correspond to
        a real MRI scan.
  \item For \textbf{each} of the four classes (Glioma, Meningioma, Pituitary Tumor,
        No Tumor), take the \emph{same} representative image you selected in
        Section~1.3 and apply your augmentation pipeline \textbf{independently four
        times}, producing four distinct augmented versions.
  \item For each class, assemble a single figure with 5 panels: the original image
        followed by its 4 augmented versions, each panel titled with the exact
        transformation(s) and parameter value(s) applied (e.g.\ \texttt{rotate
        +12\textdegree{}}, \texttt{brightness $\times$0.85}).
  \item Save the four class figures as \texttt{fig\_augmentation\_glioma.png},
        \texttt{fig\_augmentation\_meningioma.png}, \texttt{fig\_augmentation\_pituitary.png},
        and \texttt{fig\_augmentation\_notumor.png}.
\end{enumerate}
\end{stepbox}

\begin{notebox}
\textbf{Required output:} four figures total --- one per class --- each showing
1 original image plus 4 augmented versions with their transformation labels. This is
not optional and is graded as its own rubric line (see the Grading Rubric).
\end{notebox}

\begin{center}
\renewcommand{\arraystretch}{1.3}
\small\noindent\emph{Augmentation entries are ordered as flip (N/Y), rotation in degrees, zoom, brightness, contrast, and noise standard deviation.}
\begin{tabularx}{\textwidth}{|l|X|X|}
\hline
\textbf{Class} & \textbf{Representative image \newline (file/index)} & \textbf{Augmentations applied \newline (with parameters)} \\
\hline
Glioma & \texttt{00913.jpg / 912} & \footnotesize\texttt{1 N/-11.4/0.96/0.89/1.09/0.020; 2 N/-11.2/0.92/0.89/1.04/0.017; 3 Y/+2.1/1.06/0.80/1.12/0.021; 4 Y/-6.8/1.05/0.84/0.95/0.011} \\
\hline
Meningioma & \texttt{00103.png / 4269} & \footnotesize\texttt{1 Y/+4.7/0.93/0.99/1.07/0.024; 2 N/+2.6/1.00/1.03/1.05/0.028; 3 Y/-10.6/0.92/0.85/0.86/0.005; 4 Y/-2.7/1.10/0.91/0.87/0.002} \\
\hline
Pituitary Tumor & \texttt{03038.png / 13001} & \footnotesize\texttt{1 Y/+1.0/1.07/0.87/0.89/0.001; 2 N/-9.6/1.05/1.08/0.92/0.019; 3 N/+4.1/0.97/0.94/0.91/0.005; 4 N/-5.5/0.96/1.18/1.13/0.009} \\
\hline
No Tumor & \texttt{01127.jpg / 8240} & \footnotesize\texttt{1 Y/-0.3/0.92/0.94/0.83/0.000; 2 N/+11.0/1.07/1.06/1.16/0.002; 3 N/+5.7/0.94/1.14/1.09/0.002; 4 Y/-10.6/0.91/0.93/0.94/0.022} \\
\hline
\end{tabularx}
\end{center}

\figwide{fig_augmentation_glioma.png}{Glioma: original image and four augmented versions.}{fig:aug1}
\figwide{fig_augmentation_meningioma.png}{Meningioma: original image and four augmented versions.}{fig:aug2}
\figwide{fig_augmentation_pituitary.png}{Pituitary Tumor: original image and four augmented versions.}{fig:aug3}
\figwide{fig_augmentation_notumor.png}{No Tumor: original image and four augmented versions.}{fig:aug4}
\FloatBarrier

\begin{tcolorbox}[colback=lightblue!40, colframe=myColorB, title={\textbf{3.1 --- Question}}, fonttitle=\small\bfseries, breakable]
Augmentation choices that are safe for natural-image datasets (e.g.\ ImageNet) are not
automatically safe for medical imaging. For \textbf{each} augmentation you used
(horizontal flip, rotation, brightness/contrast, zoom, noise), state one reason it is
anatomically or diagnostically acceptable for brain MRI, and explain why a vertical
flip would \emph{not} be acceptable.
\end{tcolorbox}

\begin{answerbox}
Each augmentation copies a real variation. \textbf{Horizontal flip:} the brain is roughly left-right
symmetric and tumors can appear on either side, so a mirrored scan still looks like a real brain with the
same tumor type. \textbf{Rotation} ($\pm12^\circ$): the head is never placed in the scanner at exactly the
same angle. \textbf{Brightness/contrast:} different scanners and settings give different brightness.
\textbf{Zoom} (0.9--1.1): head size and field of view differ between patients, so this changes scale, not anatomy.
\textbf{Mild noise:} real MRI always has some noise.\par
\textbf{Vertical flip:} it turns the brain upside down. Real scans always have the same up-down orientation,
and position matters for diagnosis (for example, pituitary tumors sit at the base of the brain), so a vertical
flip would teach the model a wrong picture of the anatomy.
\end{answerbox}

\section{Convolution Mathematics --- Required for Every Class}

\begin{stepbox}
\begin{enumerate}[leftmargin=*, label=\arabic*.]
  \item For \textbf{each} of the four classes, take the \emph{same} representative
        image used in Sections~1.3 and~3, convert it to grayscale, and crop a small
        patch (recommended size $6 \times 6$ or $8 \times 8$ pixels) from a region that
        contains visible tumor tissue (or, for the No Tumor class, a representative
        region of healthy brain tissue). Record the pixel values of this patch.
  \item Define a single $3 \times 3$ kernel to apply to every class's patch, e.g.\ the
        vertical edge-detection kernel
        $K = \begin{bmatrix}-1 & 0 & 1\\ -1 & 0 & 1\\ -1 & 0 & 1\end{bmatrix}$.
  \item With \textbf{stride 1} and \textbf{no padding} (valid padding), manually
        compute, by hand, the output feature map values for every valid position of
        the kernel over the patch. Show the element-wise multiply-and-sum at each
        position.
  \item Reproduce the identical operation in PyTorch using \texttt{F.conv2d} on the
        same patch (as a \texttt{(1,1,H,W)} tensor) and the same kernel (as a
        \texttt{(1,1,3,3)} tensor with \texttt{bias=None}). Confirm that your manual
        values match the PyTorch output to within floating-point tolerance.
  \item Verify the output spatial dimensions using
        $\text{Output} = \left\lfloor \dfrac{\text{Input} + 2P - K}{S} \right\rfloor + 1$.
  \item Compute the receptive field of a single unit in the output of this one
        convolutional layer (equal to the kernel size here). Then compute the
        \emph{compounded} receptive field if this layer were followed by a second,
        identical $3\times3$, stride-1 convolution, using
        $\text{RF}_l = \text{RF}_{l-1} + (k_l - 1)\prod_{i<l} s_i$.
\end{enumerate}
\end{stepbox}

\begin{notebox}
\textbf{Required output:} for \textbf{each} of the four classes, one figure or
side-by-side panel showing (a) the source image with the chosen patch region
highlighted, (b) the raw patch pixel values and the kernel values, (c) your manual
feature-map calculation next to the PyTorch-computed feature map, and (d) the
receptive-field calculation for 1 and 2 stacked layers. Four complete sets are
required (Glioma, Meningioma, Pituitary Tumor, No Tumor).
\end{notebox}

\begin{center}
\renewcommand{\arraystretch}{1.3}
\footnotesize
\begin{tabularx}{\textwidth}{|l|X|c|c|c|c|c|c|}
\hline
\textbf{Class} & \textbf{Patch} & \textbf{Kernel} & \textbf{Stride} & \textbf{Pad} &
\textbf{Output} & \textbf{RF (1L)} & \textbf{RF (2L)} \\
\hline
Glioma & $6\times6$ at $(64,54)$ & $3\times3$ & 1 & 0 & $4\times4$ & $3\times3$ & $5\times5$ \\
\hline
Meningioma & $6\times6$ at $(90,38)$ & $3\times3$ & 1 & 0 & $4\times4$ & $3\times3$ & $5\times5$ \\
\hline
Pituitary Tumor & $6\times6$ at $(52,52)$ & $3\times3$ & 1 & 0 & $4\times4$ & $3\times3$ & $5\times5$ \\
\hline
No Tumor & $6\times6$ at $(61,61)$ & $3\times3$ & 1 & 0 & $4\times4$ & $3\times3$ & $5\times5$ \\
\hline
\end{tabularx}
\normalsize
\end{center}

\figwide{fig_convolution_glioma.png}{Glioma: patch, kernel, manual feature map, PyTorch feature map, and receptive field.}{fig:conv1}
\figwide{fig_convolution_meningioma.png}{Meningioma: patch, kernel, manual feature map, PyTorch feature map, and receptive field.}{fig:conv2}
\figwide{fig_convolution_pituitary.png}{Pituitary Tumor: patch, kernel, manual feature map, PyTorch feature map, and receptive field.}{fig:conv3}
\figwide{fig_convolution_notumor.png}{No Tumor: patch, kernel, manual feature map, PyTorch feature map, and receptive field.}{fig:conv4}
\FloatBarrier

\begin{tcolorbox}[colback=lightblue!40, colframe=myColorB, title={\textbf{4.1 --- Question}}, fonttitle=\small\bfseries, breakable]
\textbf{(a)} Why must your manual calculation and the PyTorch \texttt{F.conv2d} output
match exactly (within floating-point tolerance)? What would a mismatch most likely
indicate about kernel orientation (true mathematical convolution, which flips the
kernel, versus the cross-correlation that deep learning frameworks actually compute)?\\[6pt]
\textbf{(b)} Compare the receptive field after 1 layer versus after 2 stacked
$3\times3$ layers. Is this consistent with the general principle that two stacked
$3\times3$ convolutions can approximate the receptive field of one $5\times5$
convolution? Show the arithmetic.
\end{tcolorbox}

\begin{answerbox}
\textbf{(a)} \texttt{F.conv2d} does the same arithmetic as the manual loops: slide the kernel, multiply the
matching numbers, and add them up. The same operation on the same numbers must agree up to floating-point
rounding, and all four classes matched exactly (maximum absolute difference $0$). Deep learning libraries
actually compute cross-correlation (the kernel is not flipped), while true mathematical convolution flips the
kernel by $180^\circ$. For the Glioma patch, the flipped-kernel result was exactly the negative of PyTorch's
result (for example $-20$ became $20$), because this kernel flipped by $180^\circ$ is just $-K$. So a manual
result that matches PyTorch except that every sign is opposite most likely means the kernel orientation
was flipped.\par
\textbf{(b)} A $3\times3$ layer has receptive field $3$. Two stride-1 layers give
$3+(3-1)\cdot1=5$, the same as one $5\times5$ layer. The two-layer stack also uses fewer weights and has an
extra non-linearity between the layers.
\end{answerbox}

\section{Implementation and Comparison (Experiment 1)}

Implement and experimentally evaluate the following six architectures, all trained
and evaluated on the \emph{same} splits, preprocessing, and evaluation metrics from
Section~2.

\begin{itemize}[leftmargin=*]
  \item LeNet-5
  \item AlexNet
  \item A VGG-style CNN (any variant)
  \item A GoogLeNet / Inception-inspired CNN
  \item ResNet (any depth version)
  \item A lightweight EfficientNet variant
\end{itemize}

\begin{stepbox}
For \textbf{every} architecture:
\begin{enumerate}[leftmargin=*, label=\arabic*.]
  \item Implement the network in PyTorch, adapting the first convolutional layer's
        input channels to match your chosen grayscale/RGB preprocessing from Section~2.
  \item Print the complete model summary; identify which layers are convolutional,
        pooling, activation, normalization, and fully connected.
  \item Report the total number of trainable parameters.
  \item Print the output tensor shape after each major block.
  \item In one short paragraph, explain the architecture's distinguishing
        characteristic (e.g.\ inception modules, residual connections, depthwise
        separable convolutions, compound scaling).
  \item Train and evaluate with a single fixed optimizer (e.g.\ Adam), the same
        learning rate, batch size, epoch budget, and random seed used for every other
        architecture in this section. The optimizer itself is varied separately in
        Section~7.
  \item Record training time, validation performance, and test performance.
  \item If you use a pretrained backbone for any architecture, explicitly state
        whether you are doing feature extraction (frozen backbone) or fine-tuning
        (unfrozen backbone), and justify the choice.
\end{enumerate}
\end{stepbox}

\textbf{Required outputs:} training-vs-validation accuracy plot, training-vs-validation
loss plot, confusion matrix, and classification report for each architecture, plus the
consolidated comparison table below.

\begin{center}
\renewcommand{\arraystretch}{1.3}
\footnotesize\setlength{\tabcolsep}{3pt}
\begin{tabularx}{\textwidth}{|l|c|c|c|c|c|X|}
\hline
\textbf{Architecture} & \textbf{Train Acc} & \textbf{Val Acc} & \textbf{Test Acc} &
\textbf{\#Params} & \textbf{Train Time} & \textbf{Notes}\newline \tiny{(efficiency / generalization)} \\
\hline
LeNet-5 & 0.7418 & 0.7271 & 0.7374 & 1,627,916 & 215.4 s & Early stopped after 22 epochs; smallest accuracy of the six. \\
\hline
AlexNet & 0.8494 & 0.8262 & 0.8338 & 5,869,380 & 332.3 s & Completed 30 epochs; largest parameter count. \\
\hline
VGG & 0.8566 & 0.8317 & 0.8348 & 2,833,892 & 446.7 s & Early stopped after 26 epochs. \\
\hline
GoogLeNet & 0.8503 & 0.8343 & 0.8328 & 165,476 & 311.9 s & Early stopped after 23 epochs; fewest parameters. \\
\hline
ResNet & 0.8618 & 0.8322 & 0.8359 & 2,796,772 & 294.5 s & Early stopped after 21 epochs. \\
\hline
EfficientNet & 0.8638 & 0.8373 & 0.8369 & 606,276 & 504.0 s & Best test accuracy; completed 30 epochs. \\
\hline
\end{tabularx}
\end{center}

\begin{center}
\renewcommand{\arraystretch}{1.2}\footnotesize
\begin{tabular}{|l|cc|cc|cc|cc|c|}
\hline
 & \multicolumn{2}{c|}{\textbf{Glioma}} & \multicolumn{2}{c|}{\textbf{Meningioma}} & \multicolumn{2}{c|}{\textbf{Pituitary}} & \multicolumn{2}{c|}{\textbf{No Tumor}} & \\
\textbf{Architecture} & P & R & P & R & P & R & P & R & \textbf{Macro F1} \\ \hline
LeNet-5 & 0.63 & 0.76 & 0.54 & 0.35 & 0.83 & 0.93 & 0.94 & 0.89 & 0.73 \\ \hline
AlexNet & 0.77 & 0.71 & 0.68 & 0.72 & 0.92 & 0.95 & 0.98 & 1.00 & 0.84 \\ \hline
VGG & 0.76 & 0.76 & 0.69 & 0.68 & 0.93 & 0.93 & 0.98 & 0.99 & 0.84 \\ \hline
GoogLeNet & 0.77 & 0.73 & 0.67 & 0.71 & 0.92 & 0.94 & 0.99 & 1.00 & 0.84 \\ \hline
ResNet & 0.75 & 0.75 & 0.71 & 0.69 & 0.92 & 0.94 & 0.98 & 1.00 & 0.84 \\ \hline
EfficientNet & 0.77 & 0.73 & 0.69 & 0.72 & 0.92 & 0.94 & 0.99 & 1.00 & 0.84 \\ \hline
\end{tabular}

\smallskip\small Test-set classification report summary (P = precision, R = recall), CE loss with Adam.
\end{center}

\figwide{fig_architecture_comparison.png}{Architecture comparison: test accuracy, parameters, and training time.}{fig:arch}
\figpair{fig_curves_LeNet-5.png}{fig_confusion_LeNet-5.png}{LeNet-5: training vs.\ validation curves and test confusion matrix.}{fig:lenet}
\figpair{fig_curves_AlexNet.png}{fig_confusion_AlexNet.png}{AlexNet: training vs.\ validation curves and test confusion matrix.}{fig:alexnet}
\figpair{fig_curves_VGG.png}{fig_confusion_VGG.png}{VGG: training vs.\ validation curves and test confusion matrix.}{fig:vgg}
\figpair{fig_curves_GoogLeNet.png}{fig_confusion_GoogLeNet.png}{GoogLeNet: training vs.\ validation curves and test confusion matrix.}{fig:googlenet}
\figpair{fig_curves_ResNet.png}{fig_confusion_ResNet.png}{ResNet: training vs.\ validation curves and test confusion matrix.}{fig:resnet}
\figpair{fig_curves_EfficientNet.png}{fig_confusion_EfficientNet.png}{EfficientNet: training vs.\ validation curves and test confusion matrix.}{fig:effnet}
\FloatBarrier

\begin{tcolorbox}[colback=lightblue!40, colframe=myColorB, title={\textbf{5.1 --- Question}}, fonttitle=\small\bfseries, breakable]
Which architecture achieved the best test accuracy, and which was the most
computationally efficient (shortest training time or smallest parameter count)?
Were they the same architecture? Relate the differences you observed to specific
structural features (depth, residual connections, inception modules, global average
pooling, parameter sharing).
\end{tcolorbox}

\begin{answerbox}
\textbf{Best test accuracy:} EfficientNet ($0.8369$). But the top five are almost tied (AlexNet $0.8338$, VGG
$0.8348$, GoogLeNet $0.8328$, ResNet $0.8359$, EfficientNet $0.8369$): they are only $0.41$ points apart, which
is smaller than the noise of a 1,980-image test set (about $0.8$ points), so EfficientNet is not clearly better.
LeNet-5 was clearly lowest ($0.7374$).\par
\textbf{Most efficient:} GoogLeNet had the fewest parameters ($165{,}476$) and LeNet-5 had the shortest training
time ($215.4$ s, although it also stopped early at epoch 22). So they are not the same architecture, and
LeNet-5 is fast but weaker.\par
\textbf{Structure:} LeNet-5 has only 2 conv layers (fast but too simple), and about 1.61M of its 1.63M parameters
sit in the first fully connected layer. AlexNet is the biggest (5.87M, about 3.4M in two FC layers) but not better.
VGG stacks $3\times3$ convolutions and was one of the slowest ($446.7$ s). GoogLeNet uses $1\times1$ reductions
and global average pooling (its classifier has only 900 weights), which keeps it at 165k parameters. ResNet's
skip connections make a deeper network easy to train ($0.8359$ in $294.5$ s). EfficientNet reached the top
accuracy with $606{,}276$ parameters (about a tenth of AlexNet) using depthwise convolutions and
squeeze-and-excitation, but it trained the longest ($504.0$ s), probably because it has many small layers
(51 conv layers) and ran all 30 epochs.
\end{answerbox}

\section{Loss Function Comparison (Experiment 2)}

For \textbf{every} architecture from Section~5, repeat training using each of the
following loss functions, holding the dataset splits, preprocessing, optimizer,
learning rate, batch size, epochs, and seed constant:

\begin{enumerate}[leftmargin=*, label=\arabic*.]
  \item Cross-Entropy Loss.
  \item Cross-Entropy Loss with Label Smoothing (e.g.\ $\varepsilon = 0.1$).
  \item Focal Loss, implemented with the correct $(1-p_t)^{\gamma}$ modulating factor
        on the log-probability (state your chosen $\gamma$ and, given the class
        imbalance from Reflection Question~1, your chosen per-class $\alpha$ weights).
  \item Class-Weighted Cross-Entropy Loss, with per-class weights set inversely
        proportional to class frequency (state the exact weight you computed for each
        of the four classes using the counts from Reflection Question~1).
\end{enumerate}

\begin{notebox}
All four losses above are genuinely suited to this multi-class classification task.
Mean Squared Error on softmax probabilities is deliberately \textbf{excluded} from this
required comparison: its gradient saturates and shrinks when a prediction is
confidently wrong, which makes it a poor fit for a classification objective. If your
instructor wants that failure mode demonstrated directly, treat an MSE run as an
optional extension rather than part of the required comparison.
\end{notebox}

\begin{stepbox}
For each (architecture, loss) pair, record: convergence speed, training/validation
accuracy, per-class precision/recall (particularly for the minority classes), the
predicted-probability confidence distribution, and whether over/underfitting occurred.
\end{stepbox}

\begin{center}
\renewcommand{\arraystretch}{1.3}
\begin{tabular}{|l|c|c|c|c|}
\hline
\textbf{Architecture} & \textbf{Cross-Entropy} & \textbf{Label Smoothing} & \textbf{Focal Loss} & \textbf{Weighted CE} \\
\hline
LeNet-5 & 0.7374 & 0.7343 & 0.7242 & 0.7409 \\
\hline
AlexNet & 0.8338 & 0.8369 & 0.7455 & 0.8217 \\
\hline
VGG & 0.8348 & 0.8389 & 0.8177 & 0.8364 \\
\hline
GoogLeNet & 0.8328 & 0.8359 & 0.8222 & 0.8157 \\
\hline
ResNet & 0.8359 & 0.8222 & 0.8283 & 0.8364 \\
\hline
EfficientNet & 0.8369 & 0.8323 & 0.8288 & 0.8288 \\
\hline
\end{tabular}
\end{center}

\begin{center}
\renewcommand{\arraystretch}{1.2}\footnotesize
\begin{tabular}{|l|c|c|c|c|}
\hline
\textbf{Epochs to best val.\ acc.} & \textbf{CE} & \textbf{LS} & \textbf{Focal} & \textbf{WCE} \\ \hline
LeNet-5 & 17 & 13 & 14 & 12 \\ \hline
AlexNet & 28 & 30 & 13 & 26 \\ \hline
VGG & 20 & 28 & 23 & 21 \\ \hline
GoogLeNet & 19 & 17 & 13 & 7 \\ \hline
ResNet & 21 & 7 & 16 & 13 \\ \hline
EfficientNet & 30 & 13 & 30 & 21 \\ \hline
\end{tabular}
\end{center}
\begin{center}
\renewcommand{\arraystretch}{1.2}\footnotesize
\begin{tabular}{|l|cc|cc|cc|cc|}
\hline
\textbf{ResNet, loss} & \multicolumn{2}{c|}{\textbf{Glioma} (P / R)} & \multicolumn{2}{c|}{\textbf{Meningioma} (P / R)} & \multicolumn{2}{c|}{\textbf{Pituitary} (P / R)} & \multicolumn{2}{c|}{\textbf{No Tumor} (P / R)} \\ \hline
Cross-Entropy & 0.752 & 0.749 & 0.706 & 0.686 & 0.923 & 0.942 & 0.984 & 0.998 \\ \hline
Label Smoothing & 0.736 & 0.754 & 0.684 & 0.613 & 0.900 & 0.946 & 0.984 & 0.998 \\ \hline
Focal Loss & 0.736 & 0.752 & 0.684 & 0.665 & 0.930 & 0.926 & 0.998 & 0.998 \\ \hline
Weighted CE & 0.753 & 0.752 & 0.692 & 0.681 & 0.928 & 0.951 & 1.000 & 0.991 \\ \hline
\end{tabular}
\end{center}
\figwide{fig_loss_comparison.png}{Loss-function comparison: test accuracy for each architecture.}{fig:loss}
\figwide{fig_loss_confidence_resnet.png}{ResNet: predicted-probability confidence distribution for each loss function.}{fig:conf}
\FloatBarrier

\begin{tcolorbox}[colback=lightblue!40, colframe=myColorB, title={\textbf{6.1 --- Question}}, fonttitle=\small\bfseries, breakable]
\textbf{(a)} Write the gradient of cross-entropy loss with respect to the logits and
compare it, conceptually, to the gradient of focal loss. How does focal loss's
modulating term change the relative weight given to easy versus hard (frequently
misclassified) examples? \\[6pt]
\textbf{(b)} Given that Glioma is the most frequent class in this dataset, explain
whether label smoothing, focal loss, or class-weighted cross-entropy would be the most
principled choice for reducing the model's bias toward the majority class, and support
your answer with your results.
\end{tcolorbox}

\begin{answerbox}
\textbf{(a)} For softmax probabilities $p_j$ and true class $y$, cross-entropy has
$\partial L_{CE}/\partial z_j=p_j-\mathbf{1}[j=y]$. For focal loss, the gradient is the same direction multiplied
by $\alpha_y(1-p_y)^{\gamma-1}\big[(1-p_y)-\gamma p_y\ln p_y\big]$. With $\gamma=2$ (ignoring $\alpha_y$), an
easy example with $p_y=0.9$ gets a factor of only about $0.03$, while a hard example with $p_y=0.2$ gets about
$1.16$. So focal loss almost ignores easy examples and puts the learning pressure on hard,
often-misclassified ones. The weights used were $\gamma=2$ and $\alpha=w=N/(4n_c)$: Glioma $0.7917$,
Meningioma $1.1194$, Pituitary $1.0209$, No Tumor $1.1575$.\par
\textbf{(b)} Class-weighted CE is the most principled choice for the imbalance because it directly gives the
smaller classes a bigger weight. Label smoothing is not meant for imbalance, and focal loss targets hard
examples, not rare classes. But the results showed no clear benefit. Weighted CE was slightly better than CE for
LeNet-5 ($0.7409$ vs.\ $0.7374$), VGG ($0.8364$ vs.\ $0.8348$) and ResNet ($0.8364$ vs.\ $0.8359$), and worse for
AlexNet, GoogLeNet and EfficientNet. Label smoothing was higher than CE for AlexNet, VGG and GoogLeNet. For
ResNet, Meningioma recall was $0.686$ with CE, $0.681$ with weighted CE, $0.665$ with focal loss and $0.613$ with
label smoothing, and Glioma recall stayed near $0.75$ ($0.749$ to $0.754$). Focal loss gave lower accuracy than CE
for all six architectures. So in these runs no loss clearly reduced the bias toward Glioma, possibly because the
imbalance is mild ($1.46:1$) and Meningioma and Glioma are simply the hardest classes.
\end{answerbox}

\section{Optimizer Comparison (Experiment 3)}

For \textbf{every} architecture from Section~5, repeat training using each of the
following optimizers, holding the dataset splits, preprocessing, loss function
(Cross-Entropy Loss), batch size, epoch budget, and seed constant. Vary \textbf{only}
the optimizer (and its own required hyperparameters) between runs.

\begin{enumerate}[leftmargin=*, label=\arabic*.]
  \item \textbf{SGD}, momentum 0.9, Nesterov momentum enabled, learning rate 0.01.
  \item \textbf{RMSprop}, learning rate 0.001.
  \item \textbf{Adam}, learning rate 0.001.
  \item \textbf{AdamW}, learning rate 0.001, weight decay 0.01.
  \item \textbf{Nadam}, learning rate 0.001.
\end{enumerate}

\begin{notebox}
As mentioned in our discussion, \texttt{keras} supports at least 13 optimizers. However, these five optimizers were selected as the ones genuinely suited to CNN image
classification at this dataset and model scale. Optimizers built for very-large-batch
training (Lamb), extreme memory-constrained large models (Adafactor), matrix
orthogonalization in large transformer-style training (Muon), or large-scale sparse
linear models (Ftrl) are \textbf{excluded} here as they offer no demonstrated advantage,
and in some cases a documented disadvantage, for a CNN of this size on this dataset.
Treat any of them as an optional extension, not part of the required comparison.
\end{notebox}

\begin{stepbox}
For each (architecture, optimizer) pair, record: epochs to best validation accuracy
(convergence speed), final training/validation/test accuracy, total training time,
and whether the run diverged, oscillated, or plateaued early. Use a consistent
early-stopping rule (monitor validation loss, a fixed patience) across all runs so that
no optimizer is penalized or favored purely by a fixed epoch budget.
\end{stepbox}

\textbf{Required outputs:} one comparative training-loss-curve figure overlaying all
five optimizers for at least one architecture, plus the consolidated comparison table
below, \textbf{fully populated for every architecture $\times$ optimizer pair}
(6 architectures $\times$ 5 optimizers = 30 result cells).

\begin{center}
\renewcommand{\arraystretch}{1.15}
\footnotesize
\begin{tabular}{|l|c|c|c|c|c|}
\hline
\textbf{Architecture} & \textbf{SGD (Nesterov)} & \textbf{RMSprop} & \textbf{Adam} & \textbf{AdamW} & \textbf{Nadam} \\
\hline
LeNet-5      & 0.7879 & 0.7510 & 0.7374 & 0.7389 & 0.7394 \\
\hline
AlexNet      & 0.8318 & 0.7960 & 0.8338 & 0.8338 & 0.8399 \\
\hline
VGG          & 0.8424 & 0.8313 & 0.8348 & 0.8414 & 0.8369 \\
\hline
GoogLeNet    & 0.8399 & 0.8318 & 0.8328 & 0.8273 & 0.8354 \\
\hline
ResNet       & 0.8232 & 0.8399 & 0.8359 & 0.8182 & 0.8399 \\
\hline
EfficientNet & 0.8242 & 0.8308 & 0.8369 & 0.8313 & 0.8283 \\
\hline
\end{tabular}
\normalsize
\end{center}

\begin{notebox}
\textbf{Reporting convention:} record \emph{test accuracy} in the table above. In your
notebook and report, also tabulate (or log to CSV) the epochs-to-convergence and
training time for every cell, since the discussion questions below require both.
\end{notebox}

\begin{center}
\renewcommand{\arraystretch}{1.2}\footnotesize
\begin{tabular}{|l|c|c|c|c|c|}
\hline
\textbf{Epochs to best val.\ acc.} & \textbf{SGD} & \textbf{RMSprop} & \textbf{Adam} & \textbf{AdamW} & \textbf{Nadam} \\ \hline
LeNet-5 & 17 & 30 & 17 & 9 & 7 \\ \hline
AlexNet & 23 & 18 & 28 & 28 & 30 \\ \hline
VGG & 24 & 17 & 20 & 23 & 23 \\ \hline
GoogLeNet & 23 & 18 & 19 & 13 & 17 \\ \hline
ResNet & 18 & 15 & 21 & 13 & 18 \\ \hline
EfficientNet & 19 & 18 & 30 & 30 & 19 \\ \hline
\end{tabular}
\end{center}
\begin{center}
\renewcommand{\arraystretch}{1.2}\footnotesize
\begin{tabular}{|l|c|c|c|c|c|}
\hline
\textbf{Train time (s)} & \textbf{SGD} & \textbf{RMSprop} & \textbf{Adam} & \textbf{AdamW} & \textbf{Nadam} \\ \hline
LeNet-5 & 224 & 296 & 215 & 147 & 131 \\ \hline
AlexNet & 309 & 192 & 332 & 327 & 326 \\ \hline
VGG & 495 & 380 & 447 & 416 & 2,797 \\ \hline
GoogLeNet & 400 & 296 & 312 & 232 & 271 \\ \hline
ResNet & 306 & 307 & 294 & 239 & 310 \\ \hline
EfficientNet & 454 & 302 & 504 & 500 & 316 \\ \hline
\end{tabular}
\end{center}
\begin{center}
\renewcommand{\arraystretch}{1.2}\footnotesize
\begin{tabular}{|l|c|c|c|c|c|}
\hline
\textbf{Train $-$ val acc.} & \textbf{SGD} & \textbf{RMSprop} & \textbf{Adam} & \textbf{AdamW} & \textbf{Nadam} \\ \hline
LeNet-5 & 0.0177 & 0.0054 & 0.0232 & 0.0236 & 0.0190 \\ \hline
AlexNet & 0.0303 & 0.0943 & 0.0161 & 0.0154 & 0.0188 \\ \hline
VGG & 0.0256 & 0.1044 & 0.0485 & 0.0397 & 0.0158 \\ \hline
GoogLeNet & 0.0335 & 0.1543 & 0.0138 & 0.0118 & 0.0250 \\ \hline
ResNet & 0.0129 & 0.0436 & 0.0150 & 0.0177 & 0.0264 \\ \hline
EfficientNet & 0.0100 & 0.0105 & 0.0159 & 0.0067 & 0.0066 \\ \hline
\end{tabular}
\end{center}
\begin{notebox}
The VGG + Nadam training time (2,797 s) is far above every other run (131 s to 504 s). It is most likely
caused by the computer slowing down or sleeping during that run, so it is not used to compare optimizer speed.
PyTorch's default RMSprop constants ($\alpha=0.99$, $\epsilon=10^{-8}$) differ from Keras
($\rho=0.9$, $\epsilon=10^{-7}$).
\end{notebox}

\figwide{fig_optimizer_comparison.png}{Optimizer comparison: test accuracy for every architecture.}{fig:opt}
\figwide{fig_optimizer_loss_curves.png}{Training-loss curves with all five optimizers overlaid, one panel per architecture.}{fig:optcurves}
\FloatBarrier

\begin{tcolorbox}[colback=lightblue!40, colframe=myColorB, title={\textbf{7.1 --- Question}}, fonttitle=\small\bfseries, breakable]
\textbf{(a)} Across all six architectures, which optimizer most consistently achieved
the highest test accuracy? Was any optimizer's ranking architecture-dependent (i.e.\
it performed well for some architectures but poorly for others)? Cite specific rows
from your table.\\[6pt]
\textbf{(b)} SGD with Nesterov momentum typically needs more epochs to reach its best
validation accuracy than the Adam-family optimizers, but can generalize at least as
well on vision tasks once it converges. Did your results show this pattern? Support
your answer with the epochs-to-convergence and test-accuracy values you logged.\\[6pt]
\textbf{(c)} AdamW decouples weight decay from the gradient-based update, unlike Adam
with L2 regularization folded directly into the gradient. For which architecture(s) in
your table (if any) did this distinction appear to matter most? Explain why a deeper,
more over-parameterized architecture (e.g.\ VGG or ResNet) would be more sensitive to
this distinction than a shallow one (e.g.\ LeNet-5).
\end{tcolorbox}

\begin{answerbox}
\textbf{(a)} No optimizer won across all six architectures. SGD led for LeNet-5 ($0.7879$), VGG ($0.8424$) and
GoogLeNet ($0.8399$); Nadam led for AlexNet ($0.8399$); RMSprop and Nadam tied for ResNet ($0.8399$); Adam led
for EfficientNet ($0.8369$). SGD won the most rows but was the lowest for EfficientNet ($0.8242$) and second
lowest for ResNet ($0.8232$). RMSprop was tied-best for ResNet but lowest for AlexNet ($0.7960$). So the ranking
depends on the architecture.\par
\textbf{(b)} Only partly. Epochs to best validation accuracy, SGD vs.\ Adam: LeNet-5 17 vs.\ 17, AlexNet 23 vs.\ 28,
VGG 24 vs.\ 20, GoogLeNet 23 vs.\ 19, ResNet 18 vs.\ 21, EfficientNet 19 vs.\ 30. SGD was slower only for VGG and
GoogLeNet. For test accuracy, SGD was higher than Adam for LeNet-5 ($+0.0505$), VGG ($+0.0076$) and GoogLeNet
($+0.0071$), and lower for AlexNet ($-0.0020$), ResNet ($-0.0127$) and EfficientNet ($-0.0127$). So the
``slower but generalizes as well'' pattern did not show up consistently.\par
\textbf{(c)} In this notebook ``Adam'' has no weight decay, while AdamW has decoupled weight decay $0.01$. AdamW
differed most from Adam on ResNet ($0.8182$ vs.\ $0.8359$, down $0.0177$); VGG improved ($0.8348$ to $0.8414$);
LeNet-5 barely moved ($0.7374$ vs.\ $0.7389$) and AlexNet did not change. In theory, deeper networks with many
weights can overfit more, so weight decay should matter more for them than for shallow LeNet-5. The results fit
this only loosely (AdamW helped VGG but hurt ResNet), so this is not a general rule from one seed.
\end{answerbox}

\section{Kernel Size Experiment (Experiment 4)}

Using the best-performing architecture from Section~5 (or a small CNN built for this
experiment), replace its first convolutional layer's kernel size with each of:
$1\times1$, $3\times3$, $5\times5$, $7\times7$, keeping all other layers fixed.

\begin{stepbox}
\begin{enumerate}[leftmargin=*, label=\arabic*.]
  \item For each kernel size, record the resulting output feature map dimensions,
        receptive field, parameter count, training time, and test accuracy.
  \item Visualize the first-layer feature maps produced by each kernel size on the
        \emph{same} input image for at least one class.
  \item Separately compare one $5\times5$ convolution against two consecutive
        $3\times3$ convolutions: compute the effective receptive field and parameter
        count for both, and explain the nonlinearity advantage of the stacked option.
  \item Explain the role of $1\times1$ convolutions in channel transformation and
        dimensionality reduction.
\end{enumerate}
\end{stepbox}

\begin{center}
\renewcommand{\arraystretch}{1.3}
\begin{tabular}{|c|c|c|c|c|c|}
\hline
\textbf{Kernel} & \textbf{Output dims} & \textbf{Receptive field} & \textbf{\#Params} & \textbf{Train time} & \textbf{Test Acc} \\
\hline
$1\times1$ & $128\times128$ & $1 / 20$ & 23,540 total & 249.6 s & 0.6980 \\
\hline
$3\times3$ & $128\times128$ & $3 / 22$ & 23,668 total & 162.1 s & 0.7096 \\
\hline
$5\times5$ & $128\times128$ & $5 / 24$ & 23,924 total & 239.2 s & 0.7298 \\
\hline
$7\times7$ & $128\times128$ & $7 / 26$ & 24,308 total & 213.8 s & 0.7328 \\
\hline
\end{tabular}
\end{center}

\figwide{fig_kernel_comparison.png}{Kernel-size comparison (1$\times$1, 3$\times$3, 5$\times$5, 7$\times$7).}{fig:kernel}
\figwide{fig_kernel_feature_maps.png}{First-layer feature maps for each kernel size on the same Glioma image.}{fig:kernelmaps}
\begin{center}
\renewcommand{\arraystretch}{1.2}\footnotesize
\begin{tabular}{|c|c|c|c|c|c|c|}
\hline
\textbf{$C_{in}=C_{out}$} & \textbf{One 5$\times$5 params} & \textbf{RF} & \textbf{Two 3$\times$3 params} & \textbf{RF} & \textbf{Saved (\%)} & \textbf{Nonlinearities} \\ \hline
16 & 6,400 & 5 & 4,608 & 5 & 28 & 1 / 2 \\ \hline
64 & 102,400 & 5 & 73,728 & 5 & 28 & 1 / 2 \\ \hline
256 & 1,638,400 & 5 & 1,179,648 & 5 & 28 & 1 / 2 \\ \hline
\end{tabular}
\end{center}
\FloatBarrier

\begin{tcolorbox}[colback=lightblue!40, colframe=myColorB, title={\textbf{8.4 --- Question}}, fonttitle=\small\bfseries, breakable]
Based on your results, did larger kernels improve accuracy enough to justify their
extra parameters and compute cost for this dataset? Use your table's numbers to
support your answer.
\end{tcolorbox}

\begin{answerbox}
Test accuracy rose from $0.6980$ ($1\times1$) to $0.7096$ ($3\times3$), $0.7298$ ($5\times5$) and $0.7328$
($7\times7$). Total parameters only went from 23,540 to 24,308 ($+768$, about 3\%) because only the first layer
changes. First-layer compute grows with $k^2$ ($7\times7$ is about $5.4$ times the $3\times3$ layer), but training
time did not follow it (249.6, 162.1, 239.2, 213.8 s) because the runs stopped at different epochs (23, 15,
22, 19); per epoch they all took about 11 s. The gain from $3\times3$ to $5\times5$ is $+0.0202$ but from
$5\times5$ to $7\times7$ only $+0.0030$, so $5\times5$ looks like the best trade-off. This is one seed, so
differences of about 1 point should not be over-trusted.
\par\smallskip
For equal input and output channel width $C$, one $5\times5$ convolution has $25C^2$ weights, while two
$3\times3$ layers have $18C^2$, a 28\% reduction with the same $5\times5$ receptive field and an extra
non-linearity. A $1\times1$ convolution mixes channels at each pixel without changing the spatial size, so it can
change the channel count cheaply; for example, 1x1 (256 to 64) + 3x3 (64 to 64) + 1x1 (64 to 256) uses 69,632
parameters, only 11.8\% of a direct $3\times3$ convolution from 256 to 256 channels (589,824).
\end{answerbox}

\section{Stride Experiment (Experiment 5)}

Using the same base architecture as Section~8, fix the kernel size (use $3\times3$)
and vary the stride of the first convolutional layer: $S=1$, $S=2$, $S=3$.

\begin{stepbox}
\begin{enumerate}[leftmargin=*, label=\arabic*.]
  \item For each stride, verify the output dimensions using
        $\text{Output} = \lfloor (\text{Input} + 2P - K)/S \rfloor + 1$.
  \item Record the number of convolution operations performed, computational cost,
        training time, and test accuracy for each stride.
  \item Discuss how much spatial information is discarded as stride increases.
  \item Investigate whether replacing a pooling layer with a larger-stride convolution
        in the same position produces comparable downsampling and accuracy.
\end{enumerate}
\end{stepbox}

\begin{center}
\renewcommand{\arraystretch}{1.3}
\begin{tabular}{|c|c|c|c|p{6cm}|}
\hline
\textbf{Stride} & \textbf{Output dims} & \textbf{Train time} & \textbf{Test Acc} & \textbf{Notes} \\
\hline
1 & $128\times128$ & 162.1 s & 0.7096 & 40.11M conv MACs; 16,384 positions (100\%); 23,668 params. \\
\hline
2 & $64\times64$ & 283.0 s & 0.7631 & 10.03M conv MACs; 4,096 positions (25\%); 23,668 params. \\
\hline
3 & $43\times43$ & 127.1 s & 0.7500 & 4.14M conv MACs; 1,849 positions (11.3\%); 23,668 params. \\
\hline
\end{tabular}
\end{center}

\figwide{fig_stride_comparison.png}{Stride comparison: test accuracy, convolution MACs, and training time.}{fig:stride}
\FloatBarrier

\begin{tcolorbox}[colback=lightblue!40, colframe=myColorB, title={\textbf{9.2 --- Question}}, fonttitle=\small\bfseries, breakable]
Does stride affect the number of trainable parameters in a convolutional layer?
Explain why or why not, and distinguish this from stride's effect on computational
cost.
\end{tcolorbox}

\begin{answerbox}
Stride does not change a convolution's parameter count: the weights still have shape
$C_{out}\times C_{in}\times k\times k$, and all three stride rows have 23,668 parameters. It does change how many
positions the kernel visits, so the output size and compute change: the first conv output is $128\times128$
(16,384 positions), $64\times64$ (4,096) and $43\times43$ (1,849), and the total conv MACs fell from 40.11M to
10.03M to 4.14M.\par
Stride 2 keeps only 25\% of the positions and stride 3 only 11.3\%, so fine details can be skipped. Still,
accuracy did not drop: $0.7096$ (stride 1), $0.7631$ (stride 2) and $0.7500$ (stride 3), so for this small
network the lost positions did not hurt (one seed, differences near the noise level).\par
The strided-convolution downsampler gave the same sizes and reached $0.7520$ test accuracy versus $0.7096$ for
max-pooling, so it matched or beat it, but it used 72,276 parameters and 68.42M MACs versus 23,668 parameters
and 40.11M MACs.
\end{answerbox}

\section{Padding Experiment (Experiment 6)}

Compare \textbf{valid padding} ($P=0$) against \textbf{same padding} (the padding
value that keeps spatial dimensions unchanged for your chosen kernel/stride) across
the first 3 convolutional layers of your base architecture.

\begin{stepbox}
\begin{enumerate}[leftmargin=*, label=\arabic*.]
  \item Manually calculate, before running any code, the spatial dimensions after
        each of the 3 layers (L1, L2, and L3) under valid padding, and separately under same padding.
  \item Verify both sets of calculations using PyTorch.
  \item Visualize the feature maps produced at each layer under both padding schemes,
        noting information lost at the image boundary under valid padding.
  \item Record classification performance under both configurations.
  \item Explain why repeated valid-padding convolutions cause spatial resolution to
        shrink rapidly across a deep network.
\end{enumerate}
\end{stepbox}

\begin{center}
\renewcommand{\arraystretch}{1.3}
\begin{tabular}{|l|c|c|c|c|}
\hline
\textbf{Padding} & \textbf{Dims after L1} & \textbf{Dims after L2} & \textbf{Dims after L3} & \textbf{Test Acc} \\
\hline
Valid ($P=0$) & 126 & 61 & 28 & 0.7177 \\
\hline
Same & 128 & 64 & 32 & 0.7096 \\
\hline
\end{tabular}
\end{center}

\begin{notebox}
The sizes in the table are the convolution outputs. A max-pooling layer sits between the convolutions, so the full
sequence is $126\to63\to61\to30\to28\to14$ for valid padding and $128\to64\to64\to32\to32\to16$ for same padding.
\end{notebox}

\figwide{fig_padding_feature_maps.png}{Feature maps at each convolution layer under valid and same padding.}{fig:padmaps}
\figwide{fig_padding_comparison.png}{Padding comparison: dimensions and test accuracy.}{fig:pad}
\FloatBarrier

\begin{tcolorbox}[colback=lightblue!40, colframe=myColorB, title={\textbf{10.3 --- Question}}, fonttitle=\small\bfseries, breakable]
Did valid padding measurably hurt performance for a network this shallow? At what
depth (number of valid-padding layers) would the feature map shrink to a single
pixel for your chosen input size and kernel size? Show the calculation.
\end{tcolorbox}

\begin{answerbox}
Valid padding did not hurt: test accuracy was $0.7177$ versus $0.7096$ for same padding (validation $0.7064$
vs.\ $0.7089$), a small difference from one seed. The network is shallow, so losing 2 pixels per convolution at
the border costs little.\par
With a $3\times3$, stride-1 valid convolution, $n_d=128-2d$. After $d=63$ layers the map is $2\times2$, and a
64th $3\times3$ convolution does not fit, so this even input never reaches exactly $1\times1$ ($128-2d=1$ would
need $d=63.5$; an odd input such as 127 would reach it at $d=63$). For $5\times5$ kernels, $n_d=128-4d$, so 31
layers give $4\times4$.\par
Valid convolutions shrink the map quickly because the kernel cannot be centered on border pixels, so each layer
removes $k-1$ pixels from the height and from the width, and this adds up with depth (pooling shrinks it even
more). Border pixels are also used by fewer kernel windows (a corner pixel is in one $3\times3$ window, a center
pixel in nine), so information near the image edge is used less and then cut off. Same padding adds zeros so the
border pixels still get used.
\end{answerbox}

\section{Pooling Experiment (Experiment 7)}

Compare \textbf{MaxPool2D}, \textbf{AvgPool2D}, and \textbf{Global Average Pooling
(GAP)} using $2\times2$ pooling with stride 2 as the primary configuration (and
$3\times3$ pooling with stride 2 as a secondary comparison where dimensionally
appropriate).

\begin{stepbox}
\begin{enumerate}[leftmargin=*, label=\arabic*.]
  \item For at least one class's representative image, visualize the feature map
        immediately before pooling and immediately after, for both max and average
        pooling.
  \item Record output dimensions, training behavior, classification accuracy, and
        generalization performance for each pooling strategy.
  \item Explain why max pooling preserves strong local activations while average
        pooling produces smoother feature representations.
  \item Explain why pooling layers have no trainable parameters, and how gradients
        propagate differently through max pooling (routed only to the arg-max
        location) versus average pooling (distributed equally).
  \item Explain why Global Average Pooling is useful immediately before the final
        classification layer in architectures such as ResNet and GoogLeNet.
\end{enumerate}
\end{stepbox}

\begin{center}
\renewcommand{\arraystretch}{1.3}
\begin{tabular}{|p{3cm}|p{3.3cm}|c|c|p{3.6cm}|}
\hline
\textbf{Pooling} & \textbf{Output dims} & \textbf{Train time} & \textbf{Test Acc} & \textbf{Notes} \\
\hline
MaxPool $2\times2$/2 & $64\times64$ after pool 1; 16,384 features & 320.2 s & 0.7838 & Train--val gap $+0.017$; completed 30 epochs. \\
\hline
AvgPool $2\times2$/2 & $64\times64$ after pool 1; 16,384 features & 318.5 s & 0.7975 & Train--val gap $+0.008$; completed 30 epochs. \\
\hline
Global Avg Pool & $64\times64$ after pool 1; 64 features after GAP & 162.1 s & 0.7096 & Train--val gap $-0.009$; early stopped. \\
\hline
MaxPool $3\times3$/2 (secondary) & $63\times63$ after pool 1; 14,400 features & 323.2 s & 0.8000 & Train--val gap $+0.013$; completed 30 epochs. \\
\hline
AvgPool $3\times3$/2 (secondary) & $63\times63$ after pool 1; 14,400 features & 216.7 s & 0.7813 & Train--val gap $+0.016$; early stopped. \\
\hline
\end{tabular}
\end{center}

\begin{notebox}
The MaxPool and AvgPool rows use a flatten and fully connected classifier (1,072,308 parameters for $2\times2$ and
945,332 for $3\times3$), while the Global Avg Pool row is the base network with a GAP classifier (23,668 parameters).
The GAP row therefore differs from the other rows in classifier size as well as in pooling type.
\end{notebox}

\figwide{fig_pooling_feature_maps.png}{Feature map before and after max pooling and average pooling on the same trained feature map.}{fig:poolmaps}
\figwide{fig_pooling_comparison.png}{Pooling comparison: test accuracy and training time.}{fig:pool}
\FloatBarrier

\begin{tcolorbox}[colback=lightblue!40, colframe=myColorB, title={\textbf{11.2 --- Question}}, fonttitle=\small\bfseries, breakable]
In a medical imaging context, where a faint but diagnostically critical signal
(e.g.\ a small tumor boundary) could be smoothed away, which pooling strategy would
you expect to be safer, and did your results confirm this expectation?
\end{tcolorbox}

\begin{answerbox}
Max pooling keeps the biggest value in each window, so a strong activation survives unchanged, while
average pooling mixes it with weaker neighbors and gives a smoother result (in the demo, the window
$1,3,4,2$ gives a max of $4$ but an average of $2.5$). So max pooling is the reasonable expectation for a faint
localized signal. The results did not clearly confirm it: with $2\times2$ pooling, average ($0.7975$) beat max
($0.7838$), but with $3\times3$ pooling, max ($0.8000$) beat average ($0.7813$). There was no consistent
winner, and the gaps are near the noise level. Overall accuracy also cannot show whether faint tumor boundaries
were kept.\par
Pooling layers have no learned weights because they use a fixed rule. Max-pool sends the gradient only to the
winning position, and average-pool splits it equally ($0.25$ for each position of a $2\times2$ window). Global
average pooling turns each channel into one number, so the classifier is much smaller and overfits less. The GAP
row ($0.7096$) had lower accuracy than the flatten + FC rows, but it also changes the classifier size, so the
two effects cannot be separated.
\end{answerbox}

\section{Combined Hyperparameter Experiments (Experiment 8)}

After completing Sections~8--11 individually, evaluate the following six combined
configurations on your base architecture.

\begin{center}
\renewcommand{\arraystretch}{1.3}
\begin{tabular}{|c|c|c|c|c|c|}
\hline
\textbf{Config} & \textbf{Kernel} & \textbf{Stride} & \textbf{Padding} & \textbf{Pooling} & \textbf{Test Acc} \\
\hline
A & $3\times3$ & 1 & Same & Max & 0.7096 \\
\hline
B & $3\times3$ & 1 & Same & Average & 0.7444 \\
\hline
C & $5\times5$ & 1 & Valid & Max & 0.7455 \\
\hline
D & $5\times5$ & 2 & Same & Average & 0.7429 \\
\hline
E & $7\times7$ & 1 & Valid & Max & 0.7566 \\
\hline
F & $3\times3$ & 2 & Same & None & 0.7283 \\
\hline
\end{tabular}
\end{center}

\begin{center}
\renewcommand{\arraystretch}{1.2}\footnotesize
\setlength{\tabcolsep}{3pt}
\begin{tabular}{|c|c|c|c|c|c|c|p{3.9cm}|}
\hline
\textbf{Config} & \textbf{Final map} & \textbf{RF} & \textbf{Params} & \textbf{Conv MACs} & \textbf{Train time} & \textbf{Test Acc} & \textbf{Dims per layer} \\ \hline
A & $16\times16$ & 22 & 23,668 & 40.11M & 162.1 s & 0.7096 & 128, 64, 64, 32, 32, 16 \\ \hline
B & $16\times16$ & 22 & 23,668 & 40.11M & 307.3 s & 0.7444 & 128, 64, 64, 32, 32, 16 \\ \hline
C & $14\times14$ & 24 & 23,924 & 37.19M & 242.9 s & 0.7455 & 124, 62, 60, 30, 28, 14 \\ \hline
D & $8\times8$ & 43 & 23,924 & 11.08M & 131.7 s & 0.7429 & 64, 32, 32, 16, 16, 8 \\ \hline
E & $13\times13$ & 26 & 24,308 & 41.15M & 250.4 s & 0.7566 & 122, 61, 59, 29, 27, 13 \\ \hline
F & $64\times64$ & 11 & 23,668 & 94.96M & 235.2 s & 0.7283 & 64, 64, 64 \\ \hline
\end{tabular}
\end{center}

\figwide{fig_combined_comparison.png}{Combined configurations A to F: test accuracy and computation.}{fig:combined}
\FloatBarrier

\begin{tcolorbox}[colback=lightblue!40, colframe=myColorB, title={\textbf{12.1 --- Question}}, fonttitle=\small\bfseries, breakable]
Do the effects of kernel size, stride, padding, and pooling appear to be
\emph{independent}, or do you see evidence of interaction (i.e.\ the effect of one
parameter changes depending on another)? Use specific rows from your table (e.g.\
Config~C versus Config~E) to support your answer, and discuss the trade-off between
spatial resolution, receptive field, computational cost, and accuracy.
\end{tcolorbox}

\begin{answerbox}
The results suggest that the settings interact, but this six-row sweep is not a full factorial design, so it
is a hint and not proof. Changing only max to average pooling (A to B) raised accuracy from $0.7096$ to
$0.7444$ ($+0.0348$), but in Section 11 the $2\times2$ gap was only $+0.0137$ and $3\times3$ pooling reversed it, so the
pooling effect depends on the other settings. Config C ($5\times5$, valid, max, $0.7455$) vs.\ Config E ($7\times7$,
valid, max, $0.7566$) is $+0.0111$, while with same padding in Section 8 the same kernel step gave only
$+0.0030$. Config D has an $8\times8$ final map, receptive field 43 and 11.08M conv MACs, and scored $0.7429$,
almost the same as Config B ($0.7444$, 40.11M MACs). Config F keeps a $64\times64$ final map and uses the most
compute (94.96M MACs) but scored only $0.7283$. The best accuracy was Config E at a medium cost (41.15M MACs).
No configuration was best in everything: more resolution costs compute, and more downsampling gives a larger
receptive field and lower cost but a smaller final map.
\end{answerbox}

\section{Mathematical Verification Summary}

Consolidate, in your report, the following calculations (most of which you already
performed in earlier sections --- reference them rather than redoing the work):

\begin{enumerate}[leftmargin=*, label=\arabic*.]
  \item Convolution output dimensions for every configuration tested (Sections~4, 8, 9, 10).
  \item Number of convolutional parameters for each kernel size tested (Section~8).
  \item Receptive field dimensions for 1 and 2 stacked layers (Section~4) and for the
        full depth of your base architecture.
  \item Pooling output dimensions for each pooling configuration (Section~11).
  \item The flattened feature dimension entering your first fully-connected layer,
        for at least one architecture.
  \item Total trainable parameters for each of the six architectures (Section~5).
  \item The numeric loss value, computed by hand, for one specific prediction under
        at least two of the four loss functions from Section~6.
  \item The update rule for at least two of the five optimizers tested in Section~7,
        identifying which require additional per-parameter state (e.g.\ first- and
        second-moment estimates) beyond the raw gradient, and which do not.
  \item A written explanation of how gradients propagate differently through max
        pooling versus average pooling (Section~11).
\end{enumerate}

\subsection*{13.1 Output dimensions of every tested configuration}

Every size below was computed with $\lfloor (n+2P-k)/s\rfloor+1$ and matched the PyTorch output in all 15 configurations checked.
Each list shows the size after every convolution and pooling layer of the base network
(base input $128\times128$).

\begin{center}
\renewcommand{\arraystretch}{1.2}\footnotesize
\begin{tabular}{|p{8.2cm}|p{6.4cm}|}
\hline
\textbf{Configuration} & \textbf{Dimensions} \\ \hline
Kernels $1,3,5,7$ (same padding); MaxPool, AvgPool, GAP $2\times2$; Config A, B & 128, 64, 64, 32, 32, 16 \\ \hline
MaxPool, AvgPool $3\times3$/2 & 128, 63, 63, 31, 31, 15 \\ \hline
Config C ($5\times5$, valid) & 124, 62, 60, 30, 28, 14 \\ \hline
Config D ($5\times5$, stride 2) & 64, 32, 32, 16, 16, 8 \\ \hline
Config E ($7\times7$, valid) & 122, 61, 59, 29, 27, 13 \\ \hline
Config F (stride 2, no pooling) & 64, 64, 64 \\ \hline
Stride 1, 2, 3 (first layer) & 128, 64, 43 \\ \hline
Valid / same padding (Section 10) & 126, 63, 61, 30, 28, 14 / 128, 64, 64, 32, 32, 16 \\ \hline
Section 4 patch ($6\times6$, $3\times3$ kernel) & $\lfloor (6+0-3)/1\rfloor+1=4$ \\ \hline
\end{tabular}
\end{center}

\subsection*{13.2 Convolution parameters for each kernel size}

The first layer has one input channel and 16 output channels, so it has $k^2\cdot1\cdot16$ weights.
BatchNorm adds $2\cdot16=32$ parameters.

\begin{center}
\renewcommand{\arraystretch}{1.2}\footnotesize
\begin{tabular}{|c|c|c|c|}
\hline
\textbf{Kernel} & \textbf{First-layer weights} & \textbf{BatchNorm} & \textbf{Whole network} \\ \hline
$1\times1$ & 16 & 32 & 23,540 \\ \hline
$3\times3$ & 144 & 32 & 23,668 \\ \hline
$5\times5$ & 400 & 32 & 23,924 \\ \hline
$7\times7$ & 784 & 32 & 24,308 \\ \hline
\end{tabular}
\end{center}

\subsection*{13.3 Receptive fields}

One $3\times3$ layer has $\text{RF}=3$. Two stacked stride-1 layers have $3+(3-1)\cdot1=5$.
For the full depth of the base network (three convolutions with pooling in between),
the receptive field is $22\times22$. For the combined configurations it is
A 22, B 22, C 24, D 43, E 26, and F 11.

\subsection*{13.4 Pooling output dimensions}

\begin{center}
\renewcommand{\arraystretch}{1.2}\footnotesize
\begin{tabular}{|c|c|c|c|}
\hline
\textbf{Input} & \textbf{Pool / stride} & \textbf{$\lfloor (n-k)/s\rfloor+1$} & \textbf{PyTorch} \\ \hline
128 & 2 / 2 & 64 & 64 \\ \hline
128 & 3 / 2 & 63 & 63 \\ \hline
126 & 2 / 2 & 63 & 63 \\ \hline
126 & 3 / 2 & 62 & 62 \\ \hline
128 & global & 1 & $1\times1$ \\ \hline
\end{tabular}
\end{center}

\subsection*{13.5 Flattened feature dimension before the first fully connected layer}

For LeNet-5 by hand: $128\to124$ (first $5\times5$ convolution) $\to62$ (average pooling)
$\to58$ (second $5\times5$ convolution) $\to29$ (average pooling). With 16 channels,
$16\cdot29\cdot29=13{,}456$ features, which matches the model. The same value for every architecture:

\begin{center}
\renewcommand{\arraystretch}{1.2}\footnotesize
\begin{tabular}{|l|c|c|c|c|c|c|}
\hline
\textbf{Architecture} & LeNet-5 & AlexNet & VGG & GoogLeNet & ResNet & EfficientNet \\ \hline
\textbf{Features into classifier} & 13,456 & 2,304 & 1,024 & 224 & 256 & 320 \\ \hline
\end{tabular}
\end{center}

\subsection*{13.6 Total trainable parameters}

LeNet-5 1,627,916; AlexNet 5,869,380; VGG 2,833,892; GoogLeNet 165,476; ResNet 2,796,772;
EfficientNet 606,276 (Section~5 table).

\subsection*{13.7 Loss values computed by hand for one prediction}

The logits are $z=(2.0,\,0.5,\,-1.0,\,0.0)$ and the true class is Meningioma (index 1).
The softmax is $(0.7101,\,0.1584,\,0.0354,\,0.0961)$, so $p_y=0.1584$.

\begin{center}
\renewcommand{\arraystretch}{1.3}\footnotesize
\begin{tabular}{|l|l|c|}
\hline
\textbf{Loss} & \textbf{Calculation} & \textbf{Value} \\ \hline
Cross-entropy & $-\ln(0.1584)$ & 1.842350 \\ \hline
Label smoothing ($\varepsilon=0.1$) & $0.9(1.8423)+0.1\cdot\text{mean}(-\ln p_c)$ & 1.854850 \\ \hline
Focal ($\gamma=2$, $\alpha_y=1.1194$) & $-1.1194\,(1-0.1584)^2\ln(0.1584)$ & 1.460525 \\ \hline
Weighted CE (before batch reduction) & $1.1194\times1.8423$ & 2.062262 \\ \hline
\end{tabular}
\end{center}

The first three values matched the PyTorch loss functions. PyTorch's weighted CE for a single sample is
1.842350 because the weight cancels in the weighted mean.

\subsection*{13.8 Optimizer update rules and per-parameter state}

\textbf{SGD with Nesterov momentum:} $v_t=\mu v_{t-1}+g_t,\quad \theta_{t+1}=\theta_t-\eta\,(g_t+\mu v_t)$.\\
\textbf{Adam:} $m_t=\beta_1 m_{t-1}+(1-\beta_1)g_t,\quad v_t=\beta_2 v_{t-1}+(1-\beta_2)g_t^2,\quad
\theta_{t+1}=\theta_t-\eta\,\dfrac{\hat m_t}{\sqrt{\hat v_t}+\epsilon}$, where
$\hat m_t=m_t/(1-\beta_1^t)$ and $\hat v_t=v_t/(1-\beta_2^t)$.

For gradients $(0.5,-0.2,0.3)$, $\theta_0=1$, and $\eta=0.01$, the manual updates gave
0.9905, 0.99025, and 0.982525 for SGD-Nesterov
and 0.99, 0.98654, and 0.98110 for Adam. Both matched \texttt{torch.optim} exactly.

\begin{center}
\renewcommand{\arraystretch}{1.2}\footnotesize
\begin{tabular}{|l|c|l|}
\hline
\textbf{Optimizer} & \textbf{State tensors} & \textbf{State names} \\ \hline
SGD (no momentum) & 0 & none (only the raw gradient) \\ \hline
SGD with momentum & 1 & momentum buffer \\ \hline
RMSprop & 1 & square\_avg \\ \hline
Adam, AdamW, Nadam & 2 & exp\_avg, exp\_avg\_sq \\ \hline
\end{tabular}
\end{center}

\subsection*{13.9 Gradient flow through max pooling versus average pooling}

\figwide{fig_pooling_gradient_flow.png}{Gradient with respect to the input for a $4\times4$ example with an upstream gradient of 1.}{fig:gradflow}

Max pooling sends the whole gradient to the one position that held the maximum, and the other three
positions in each window get 0. Average pooling splits the gradient equally, so every position in a
$2\times2$ window gets $0.25$. Pooling layers have no trainable parameters because they apply a fixed rule
(maximum or mean) with no weights to learn.
\FloatBarrier

\section{Visualization Requirements --- Checklist}

Your notebook and report must include all of the following, each with a title, axis
labels, and (where applicable) a legend:

\begin{itemize}[leftmargin=*]
  \item Original input images for all four classes (Section~1).
  \item Per-class data augmentation figures --- 4 total (Section~3).
  \item Per-class convolution calculation figures --- 4 total, each showing patch,
        kernel, manual feature map, and PyTorch-verified feature map (Section~4).
  \item Intermediate convolutional feature maps for the kernel-size experiment (Section~8).
  \item Feature maps before and after padding, for both padding schemes (Section~10).
  \item Feature maps before and after max pooling and average pooling (Section~11).
  \item Training-vs-validation loss curves and accuracy curves for every architecture (Section~5).
  \item Confusion matrices for every architecture (Section~5).
  \item Architecture comparison chart (Section~5), loss-function comparison
        chart (Section~6), and optimizer comparison chart (Section~7).
  \item Kernel-size, stride, and padding comparison charts (Sections~8--10).
\end{itemize}

\section{Implementation and Code Organization Requirements}

\begin{stepbox}
Organize your source code into the following reusable modules rather than one script:
\end{stepbox}

\begin{lstlisting}[style=pythonstyle]
project/
    dataset.py      # ImageFolder loading, stratified split, augmentation pipeline
    models.py       # LeNet5, AlexNet, VGG, GoogLeNet, ResNet, EfficientNet classes
    losses.py       # CrossEntropy, LabelSmoothingCE, FocalLoss, WeightedCE
    train.py         # training/validation loop, optimizer factory, checkpointing
    evaluate.py     # test-set evaluation, confusion matrix, classification report
    experiments.py  # kernel/stride/padding/pooling sweep configurations
    visualize.py    # feature map, augmentation, and curve plotting utilities
    main.py          # orchestrates everything via configuration dictionaries
\end{lstlisting}

Every convolutional, pooling, and loss configuration must be changeable through a
configuration dictionary or function arguments, \textbf{without} rewriting the model
class. Include model checkpointing (save the best validation-accuracy weights per
architecture) and basic experiment logging (a CSV or dictionary of every run's
settings and results).

\begin{notebox}
\textbf{Code organization:} the project folder contains the eight modules listed above, plus \texttt{config.py}
(settings, seed, paths) and \texttt{math\_checks.py} (hand calculations checked against PyTorch).
\texttt{main.py} runs any experiment from one command, for example \texttt{python main.py --stage loss}.
Every convolution, pooling, and loss setting is changed through a configuration dictionary or function arguments
(\texttt{flex\_cfg}, \texttt{make\_spec}), so no model class has to be rewritten. Each run saves its
best-validation weights in \texttt{checkpoints/}, its results in \texttt{runs/}, and a row in
\texttt{experiment\_log.csv}. Every run uses \texttt{SEED = 42}, batch size 32, a maximum of 30 epochs, and
patience 6. The modules contain the same code as the notebook and produce the same run identifiers, so the
finished runs are reused and not trained again.
\end{notebox}

\section*{Technical Report}

\textbf{Objective and setup.} The goal was to classify brain MRI images into four classes (Glioma, Meningioma,
Pituitary Tumor, No Tumor) and to measure how architecture, loss function, optimizer, kernel size, stride,
padding, and pooling affect performance. All 13,196 images were converted to $128\times128$ grayscale, split
70/15/15 with a stratified split, and normalized with training-split statistics. Every model was trained from
scratch on a GTX 1080 with batch size 32, a maximum of 30 epochs, early stopping on validation loss (patience 6),
and \texttt{SEED = 42}. Sections 5 to 7 compare six architectures. Sections 8 to 12 use one small CNN with 23,668
parameters, so their accuracies (0.70 to 0.80) are not comparable with the six architectures (0.83 to 0.84).

\textbf{Main results.}

\begin{center}
\renewcommand{\arraystretch}{1.2}\footnotesize
\begin{tabularx}{\textwidth}{|l|X|c|}
\hline
\textbf{Experiment} & \textbf{Best observed result} & \textbf{Test Acc} \\ \hline
Architectures (CE, Adam) & EfficientNet & 0.8369 \\ \hline
Loss functions & VGG with label smoothing & 0.8389 \\ \hline
Optimizers (CE) & VGG with SGD (Nesterov) & 0.8424 \\ \hline
Kernel size & $7\times7$ first layer & 0.7328 \\ \hline
Stride & Stride 2 (about 25\% of the conv positions) & 0.7631 \\ \hline
Padding & Valid & 0.7177 \\ \hline
Pooling & MaxPool $3\times3$/2 (secondary); AvgPool $2\times2$ is the best primary setting (0.7975) & 0.8000 \\ \hline
Combined & Config E ($7\times7$, stride 1, valid, max) & 0.7566 \\ \hline
\end{tabularx}
\end{center}

\textbf{Discussion.} Observed: the five larger architectures finished within 0.4 percentage points of each other
under the base setting (0.8328 to 0.8369), while LeNet-5 was clearly lower (0.7374). A test set of 1,980 images has
a sampling error of about 0.8 percentage points at 84\% accuracy, so the differences among the five larger models are
within noise. Class-wise, No Tumor was almost perfect (recall 0.99 to 1.00), Pituitary Tumor reached 0.93 to 0.95,
and Meningioma was the hardest class (recall 0.68 to 0.72, except LeNet-5 at 0.35), followed by Glioma
(0.71 to 0.76). No optimizer or loss function won across all architectures. SGD improved LeNet-5 the most
($0.7374\to0.7879$). Focal loss lowered accuracy for every architecture, and inverse-frequency weighting did not
raise Meningioma recall for ResNet. A stride of 2 cut the convolution cost to a quarter and still scored higher
than stride 1 (0.7631 vs.\ 0.7096). Expected from theory but not confirmed: SGD was not consistently slower than the
Adam family, and max pooling did not beat average pooling at $2\times2$ (0.7838 vs.\ 0.7975).

\textbf{Limitations.} Every result comes from a single seed, so differences below about one percentage point should
not be treated as real. The pooling comparison changes the classifier head together with the pooling type
(Section~11), and the VGG + Nadam training time is an outlier. An accuracy near 84\% is not good enough for
clinical use.

\textbf{Conclusion.} The best single result in the six-architecture comparison was VGG with SGD (Nesterov) at
0.8424. GoogLeNet gave almost the same accuracy (0.8328 with Adam, 0.8399 with SGD) with only 165,476 parameters,
which makes it the smallest model.

\section*{Executive Summary}

\begin{itemize}[leftmargin=*]
  \item Six CNN architectures classified four brain MRI classes with 0.74 to 0.84 test accuracy, and every run
        used the same split, preprocessing, and seed.
  \item The most effective configuration was VGG with cross-entropy and SGD (Nesterov), at 0.8424 test accuracy.
        The smallest was GoogLeNet, with 165,476 parameters and 0.8328 to 0.8399 accuracy.
  \item The top five architectures differ by less than the sampling error of the test set, so none can be called
        clearly better.
  \item Meningioma and Glioma are the hardest classes; No Tumor and Pituitary Tumor are easy.
  \item For the five larger architectures, loss function and optimizer choices moved test accuracy by about 1 to 2
        points. The largest changes were LeNet-5 with SGD (+0.0505 over Adam), AlexNet with focal loss
        ($-0.0883$ versus CE), and AlexNet with RMSprop ($-0.0378$ versus Adam).
  \item In the stride experiment, stride 2 had the best accuracy (0.7631) with a quarter of the convolution positions.
  \item These are single-seed results, so they show what happened in these runs and do not prove general rules.
\end{itemize}

\section*{Submission Checklist}

Before submitting, verify that you have:

\begin{itemize}[leftmargin=2em, label=$\square$]
  \item Completed all coding cells in your own notebook, following the given instructions.
  \item Produced the \textbf{4 per-class data augmentation figures} (Section~3).
  \item Produced the \textbf{4 per-class convolution calculation figures}, with manual
        values matching PyTorch output (Section~4).
  \item Filled in every results table in Sections~5--12 with your actual values.
  \item Answered every numbered discussion question (Sections~1--12) in this activity sheet.
  \item Attached or embedded all required figures (Section~14).
  \item Organized your code according to Section~15 and included model checkpoints for
        at least your best-performing architecture per experiment.
  \item Written a concise technical report and an executive summary identifying the
        most effective configuration(s), clearly distinguishing observed results from
        theoretical expectations.
  \item Submitted (a) this completed activity sheet, (b) your \texttt{.ipynb} notebook
        and supporting \texttt{.py} modules, and (c) your technical report.
\end{itemize}

\section*{Grading Rubric}

\begin{center}
\renewcommand{\arraystretch}{1.4}
\begin{tabular}{|p{10cm}|c|}
\hline
\textbf{Assessment Component} & \textbf{Weight} \\
\hline
CNN architecture implementation and comparison\newline (Section~5) & 15 \\
\hline
Loss function implementation and comparison\newline  (Section~6) & 10 \\
\hline
Optimizer implementation and comparison\newline  (Section~7) & 10 \\
\hline
Kernel size, stride, and padding experiments\newline  (Sections~8--10) & 12 \\
\hline
Max pooling, average pooling, and GAP experiments\newline  (Section~11) & 8 \\
\hline
Combined hyperparameter experiments\newline  (Section~12) & 5 \\
\hline
Per-class data augmentation demonstration \newline (Section~3) & 5 \\
\hline
Per-class convolution calculation \& verification \newline (Section~4) & 8 \\
\hline
Mathematical verification summary \newline (Section~13) & 8 \\
\hline
Visualization and experimental analysis \newline (Section~14) & 9 \\
\hline
Code quality, reproducibility, and documentation \newline (Section~15) & 5 \\
\hline
Technical report and interpretation & 5 \\
\hline
\textbf{Total} & \textbf{100} \\
\hline
\end{tabular}
\end{center}

\end{document}
